% Part: proof-theory
% Chapter: natural-deduction
% Section: rules-proofs

\documentclass[../../../include/open-logic-section]{subfiles}

\begin{document}

\olfileid{pt}{nat}{rp}

\olsection{Rules and \usetoken{P}{proof}}

We begin by giving some definitions and fixing some terminology.

Consider the two rules in \Log{N1ci} that mention $!A
\lif !B$.
\[
\bottomAlignProof
\AxiomC{$\Discharge{!A}{x}$}
\shortDeduce
\DeduceC{$!B$}
\DischargeRule{\Intro{\lif}}{x}
\UnaryInfC{$!A \lif !B$}
\DisplayProof
\qquad
\bottomAlignProof
\AxiomC{$!A \lif !B$}
\AxiomC{$!A$}
\RightLabel{\Elim{\lif}}
\BinaryInfC{$!B$}
\DisplayProof
\]
The \Elim{\lif} rule is simple enough. The !!{formula}s above the line
are the \emph{premises}. The one on the left the \emph{principal}
!!{formula}~$!A \lif !B$. All !!{elimination} rules have one such
premise, always on the far left. It is called the \emph{major} premise
of the inference. If there are others, such as an this case, they are
called \emph{minor} premises. The premise(s) of !!{introduction} rules
are both called major premise(s). The !!{formula} below the line is the
\emph{conclusion}.

The \Intro{\lif} rule only has one premise, and here the conclusion
is the principal formula~$!A \lif !B$. All !!{introduction} rules are
like this: the conclusion is the principal formula; its main operator
is the one being ``introduced.''

The notation $\Discharge{!A}{x}$ means the following. In !!a{proof},
the \Intro{\lif} inference is applied to !!a{formula}~$!B$ as premise.
The sub-!!{proof} ending in~$!B$ has a number of initial !!{formula}s.
These are called \emph{assumptions}. !!^a{proof} of~$!B$ from one or
more assumptions of the form~$!A$ is !!a{proof} that $!A$
entails~$!B$. When we apply the \Intro{\lif} rule, we conclude~$!A
\lif !B$, and the conclusion no longer depends on the assumption~$!A$.
To indicate this, the assumptions of the form~$!A$ in the !!{proof}
including the~\Intro{\lif} inference are \emph{discharged} or
\emph{closed}. The rule states which assumptions can be discharged. In
the case of \Intro{\lif} with conclusion $!A \lif !B$, that's
assumptions that have the form~$!A$, i.e., agree with the antecedent
of the conditional. It is sometimes important to indicate which
assumptions are discharged at which rule, and this is accomplished
with the discharge label~$x$. Importantly, rules that discharge
assumptions don't \emph{have to} discharge \emph{all} assumptions of the
requisite form, or indeed any at all. For instance, the following is
correct even though the second \Intro{\lif} does not discharge any
assumptions:
\[
\AxiomC{$\Discharge{!A}{y}$}
\DischargeRule{\Intro{\lif}}{x}
\UnaryInfC{$!B \lif !A$}
\DischargeRule{\Intro{\lif}}{y}
\UnaryInfC{$!A \lif (!B \lif !A)$}
\DisplayProof
\]
When a rule does not discharge any assumptions, we just leave out the
discharge label (as we did here).

In the !!{proof}
\[
\AxiomC{$\Discharge{!A}{x}$}
\AxiomC{$\Discharge{!A}{y}$}
\RightLabel{\Intro{\land}}
\BinaryInfC{$!A \land !A$}
\RightLabel{\Elim{\land}}
\UnaryInfC{$!A$}
\DischargeRule{\Intro{\lif}}{x}
\UnaryInfC{$!A \lif !A$}
\DischargeRule{\Intro{\lif}}{y}
\UnaryInfC{$!A \lif (!A \lif !A)$}
\DisplayProof
\]
the first \Intro{\lif} inference discharges the left assumption
$!A^x$, the second the right assumption~$!A^y$. That is, all
assumptions labelled~$x$ are discharged by the \Intro{\lif} with
label~$x$, and all assumptions with label~$y$ by the one labelled~$y$.
For this to work, it is important that all assumptions with the same
label are also the same !!{formula}.

The quantifier rules are slightly more complicated. E.g., the rules in
\Log{N1i} for $\lexists$ are:
\[
\bottomAlignProof
\AxiomC{$!A(t)$}
\RightLabel{\Intro{\lexists}}
\UnaryInfC{$\lexists[x][!A(x)]$}
\DisplayProof
\qquad
\bottomAlignProof
\AxiomC{$\lexists[x][!A(x)]$} 
\AxiomC{$\Discharge{!A(c)}{x}$}
\shortDeduce
\DeduceC{$!B$}
\DischargeRule{\Elim{\lexists}}{x}
\BinaryInfC{$!B$} \DisplayProof 
\bottomAlignProof
\] In the
\Intro{\lexists} rule, the premise is~$!A(t)$ where $t$~is a closed
term, i.e., a term without variables. Such an inference is
correct---that is, it matches the schematic rule \Intro{\lexists}---if
there is some !!{formula}~$!A$, some variable~$x$, and some closed
term~$t$ such that the conclusion is $\lexists[x][!A]$ and premise
is~$\Subst{!A}{t}{x}$. The \Elim{\lexists} allows us to discharge
assumptions of the form~$!A(c)$, i.e., for every inference there is
some !!{constant}~$c$ and assumptions of the form $\Subst{!A}{c}{x}$
can be discharged. The discharged assumptions of one inference all
have to use the same~$c$. This constant is called the
\emph{eigenvariable} of the \Elim{\lexists} inference. The inference
is only correct if $c$~does not occur in an open assumption other than
the discharged~$!A(c)$, does not occur in~$!B$, and does not occur
in~$!A$. This is called the \emph{eigenvariable condition}.

!!^{proof}s in \Log{N1ci} are finite trees of sequents that are
inductively generated from axioms using the rules of inference. Each
leaf is an assumption, possibly labelled with a discharge label, and
every other !!{formula} follows from the !!{formula}s immediately
above it by one of the inference rules of the system. If an occurence
of a !!{formula} as an assumption in the tree above an inference has
not been discharged by an inference above it, it is called \emph{open}
(at the inference).

\begin{defn}\ollabel{defn:proof}
\emph{!!^{proof}s} in \Log{N1ci} are finite trees of
!!{formula}s inductively as follows:
\begin{enumerate}
  \item\ollabel{defn:prf:assumption} Every !!{formula} 
  labelled by a discharge label such as $x$ or $y$, by itself is
  !!a{proof}. In !!a{proof} consisting of a single !!{formula}, the
  !!{formula} counts as an open assumption.
  \item\ollabel{defn:prf:one-prem} If $R$ is a rule with one, two, or
  three premises $!A_1$, $!A_2$, $!A_3$, conclusion~$!A$, and
  $\delta_1$, $\delta_2$, $\delta_3$ are !!{proof}s ending in $!A_1$,
  $!A_2$, $!A_3$, respectively, the
  \[
  \AxiomC{}
  \RightLabel{$\delta_1$}
  \DeduceC{$!A_1$}
  \RightLabel{$R$}
  \UnaryInfC{$!A$}
  \DisplayProof
\qquad
  \AxiomC{}
  \RightLabel{$\delta_1$}
  \DeduceC{$!A_1$}
  \AxiomC{}
  \RightLabel{$\delta_2$}
  \DeduceC{$!A_2$}
  \RightLabel{$R$}
  \BinaryInfC{$!A$}
  \DisplayProof
\qquad
  \AxiomC{}
  \RightLabel{$\delta_1$}
  \DeduceC{$!A_1$}
  \AxiomC{}
  \RightLabel{$\delta_2$}
  \DeduceC{$!A_2$}
  \AxiomC{}
  \RightLabel{$\delta_3$}
  \DeduceC{$!A_2$}
  \RightLabel{$R$}
  \TrinaryInfC{$!A$}
  \DisplayProof
  \]
  is !!a{proof}, respectively, provided the assumption labels in the
  !!{proof}s are compatible (no two different !!{formula}s have the
  same discharge label) and the inferences are correct applications
  of~$R$.
  \item Topmost !!{formula} occurrences in the new !!{proof} count as
  \emph{open assumptions} of the !!{proof} if they are open
  assumptions of $\delta_1$, $\delta_2$, or~$\delta_3$, except those
  occurrences discharged by the rule. If the rule is labelled~$x$ (or
  $x\, y$), then for \Intro{\lif} and~\Intro{\lnot}, that is all open
  assumptions in~$\delta_1$ labelled~$x$; for \Elim{\lexists}, it is
  all open assumptions labelled~$x$ in~$\delta_2$; for \Elim{\lor}, it
  is all open assumptions labelled~$x$ in~$\delta_2$ and all open
  assumptions labelled~$y$ in~$\delta_3$.
\end{enumerate}
\end{defn}

The !!{proof}s used in the inductive construction of !!a{proof} are
called its \emph{sub-!!{proof}s}. The sub-!!{proof}s ending in the
premises of an inference are called the \emph{immediate
sub-!!{proof}s} of the inference. The rules of \Log{N2c} and \Log{N2i}
are given in \olref{tab:N1}

\subfile{rules-N1}

\begin{defn}
The \emph{!!{height}}~$\pheight{\delta}$ of !!a{proof}~$\delta$ is
inductively defined as follows.
\begin{enumerate}
  \item If $\delta$ consists of an assumption, $\pheight{\delta} = 0$.
  \item If $\delta$ ends in an inference with one premise with immediate
  sub-!!{proof}~$\delta_1$, then $\pheight{\delta} = \pheight{\delta_1} + 1$.
  \item If $\delta$ ends in an inference with two premises with immediate
  sub-!!{proof}s~$\delta_1$ and~$\delta_2$, then $\pheight{\delta} =
  \max(\pheight{\delta_1}, \pheight{\delta_2}) + 1$.
  \item If $\delta$ ends in \Elim{\lor} with immediate
  sub-!!{proof}s~$\delta_1$, $\delta_2$, and~$\delta_3$, then $\pheight{\delta} =
  \max(\pheight{\delta_1}, \pheight{\delta_2}, \pheight{\delta_3}) + 1$.
\end{enumerate}
If a sequent~$S$ has !!a{proof} of height $\le n$, we write $\Proves[n] S$.
\end{defn}

\begin{ex}
In \Log{N1i}, any !!{formula} counts as !!a{proof} from itself as an
open assumption. So
\[
!A \land !B^x
\]
is !!a{proof}~$\delta_1$. Its height is~$0$.
From $\delta_1$ we can obtain two !!{proof}s~$\delta_2$ and~$\delta_3$
by applying one of the two version of~\Elim{\land}
\[
\AxiomC{$!A \land !B^x$}
\RightLabel{\Elim{\land}[2]}
\UnaryInfC{$!B$}
\DisplayProof
\qquad
\AxiomC{$!A \land !B^x$}
\RightLabel{\Elim{\land}[1]}
\UnaryInfC{$!A$}
\DisplayProof
\]
The immediate sub-!!{proof}s of the last inference in~$\delta_2$ and
in~$\delta_3$ is the same: $!A \land !B$. Both !!{proof}s have the
occurrence of $!A\land !B$ as an open assumption. We have
$\pheight{\delta_1} = \pheight{\delta_2} = 1$.

The rule \Intro{\land} requires two premises of the form $!B$ and $!A$
to justify a conclusion of the form~$!B \land !A$. Thus, the following
is !!a{proof}~$\delta_4$:
\[
\AxiomC{$!A \land !B^x$}
\RightLabel{\Elim{\land}[2]}
\UnaryInfC{$!B$}
\LeftSubproofLabel{$\delta_2$}
\AxiomC{$!A \land !B^x$}
\RightLabel{\Elim{\land}[1]}
\UnaryInfC{$!A$}
\RightSubproofLabel{$\delta_3$}
\RightLabel{\Intro{\land}}
\BinaryInfC{$!B \land !A$}
\DisplayProof
\]
Its height $\pheight{\delta_4} = \max(\pheight{\delta_2},
\pheight{\delta_3})+1 = \max(1,1) + 1 = 2$. It has two occurences of
$!A \land !B$ as assumptions, both are open.

Finally we can apply \Intro{\lif} to obtain $\delta_5$:
\[
\AxiomC{$\Discharge{!A \land !B}{x}$}
\RightLabel{\Elim{\land}[2]}
\UnaryInfC{$!B$}
\LeftSubproofLabel{$\delta_2$}
\AxiomC{$\Discharge{!A \land !B}{x}$}
\RightLabel{\Elim{\land}[1]}
\UnaryInfC{$!A$}
\RightSubproofLabel{$\delta_3$}
\RightLabel{\Intro{\land}}
\BinaryInfC{$!B \land !A$}
\DischargeRule{\Intro{\lif}}{x}
\UnaryInfC{$(!A \land !B) \lif (!B \land !A)$}
\DisplayProof
\]
Its height is $\pheight{\delta_4} + 1 = 3$. The \Intro{\lif} inference
discharges all occurrences of assumptions of the form $!A \land !B$
labelled~$x$, i.e., both previously open assumptions of the immediate
subproof. Since those are the only assumptions of $\delta_5$, it has
no open assumptions, and so is !!a{proof} \emph{of} its end-!!{formula}
$(!A \land !B) \lif (!B \land !A)$.
\end{ex}

\end{document}